\section{Plataforma experimental y calibración}
\subsection{Instrumentación}
La plataforma comprende un PSoC 5LP con ADC, DAC y temporizador de muestreo;
planta analógica con acondicionamiento de entrada y salida; fuente bipolar
Minipa MPL-3305M ajustada a $\pm\SI{10.2}{\volt}$; PC con MATLAB; y
osciloscopio \emph{Rigol DS1202Z-E}. Se emplearon sondas $10\times$,
acoplamiento DC, masas comunes, telemetría UART a
\SI{57600}{bit\per\second} y adquisición sincronizada. J5 midió el mando
físicamente aplicado y J7 la salida de la planta.

\subsection{Calibración y revalidación}
La Tabla~\ref{tab:calib_historica} reúne los parámetros obtenidos en el
Laboratorio 5 y reutilizados como punto de partida. Antes de cerrar el lazo se
realizó un ensayo abierto identificado de aproximadamente
\SI{50}{\milli\volt}. La variación robusta medida fue
$\Delta J5=\SI{45.38}{\milli\volt}$ y
$\Delta J7=\SI{44.45}{\milli\volt}$, equivalente a una ganancia incremental
de 0,98. La concordancia entre ambos canales confirmó polaridad, continuidad,
masa común y rango de señal antes de habilitar los controladores.

\begin{table}[t]
\centering
\caption{Parámetros de calibración adoptados del Laboratorio 5.}
\label{tab:calib_historica}
\small
\begin{tabular}{@{}ll@{}}
\toprule
Parámetro & Valor documentado \\
\midrule
Offset del ADC & \SI{2.353802067}{\volt} \\
Escala ADC/planta & \SI{0.210251951}{\volt\per\volt} \\
Ganancia de DAC (+) & 3,842612 \\
Ganancia de DAC (-) & 2,640105691 \\
Offset físico J5 & \SI{17.209}{\milli\volt} \\
Factor de escala J5/mando aplicado & 0,981710 \\
RMSE del ajuste de J7 & \SI{3.374}{\milli\volt} \\
\bottomrule
\end{tabular}
\end{table}

En cada ensayo se conservaron las ondas crudas de ambos canales, la telemetría
digital, la referencia, los estados estimados, el mando solicitado, el mando
nominal del DAC, banderas, tiempo de ejecución y metadatos de disparo. Esta
trazabilidad permite comparar resultados calculados, simulados y experimentales
sin inferir tensiones físicas únicamente a partir de códigos digitales.
